\section{Introduction}\label[section]{sec:intro}

When cost accrues throughout processing, it depends on when each fraction of a job is served. We study a positively weighted sum of first processing moments, where a job's moment is the integral of time over the set on which it runs. The machines are identical and allow preemption and migration, and all jobs are available at time zero. A job may occupy at most one machine at a time.

Queyranne posed the complexity question for this model in 2010 \cite[pp.~15--17]{queyranne2010}. The note describes a fixed-machine convex quadratic formulation with a variable for each feasible job subset. The number of subsets of size at most $m$ is polynomial in $n$ for fixed $m$, but can be exponential in $n$ when $m$ varies. An algorithm with a polynomial binary-time bound for a variable machine count must avoid that enumeration.

We prove such a bound and give an algorithm returning rational schedule records and a rational certificate of the exact optimum. We first relax binary activities to bounded fractional rates. Sorting score zeros and crossings allows exact integration of the support function of the rate polytope. A linear feasibility problem tests dual optimality even at tied scores. A description bound for the polytope of dual minimizers then enables exact strong optimization in the framework of Gr\"otschel, Lov\'asz and Schrijver \cite{gls1988}. The description bound uses a rational quadratic representation on each parameter cell. The algorithm needs neither the full arrangement nor a time discretization. A second result gives an explicit finite schedule from the optimal rates. Its reflection step preserves the first moment of each job, as well as its work.

\Cref{sec:model} specifies the model and results. Attained duality, the minimizer polytope, exact recovery and unrestricted optimality are developed in \Cref{sec:dual,sec:separation,sec:recovery}. The realization and the exact certificate for Queyranne's example are in \Cref{sec:realization,sec:example}. The explicit bit bounds and imported optimization contracts appear in \Cref{app:encoding,app:gls}. \Cref{sec:formal} explains the scope of the Lean proofs.

\section{Related work}\label[section]{sec:related}

\paragraph{Scheduling.}
Goemans \cite{goemans1997} used mean busy times in approximation algorithms for scheduling with release dates. Subsequent work analyzed linear relaxations through this objective \cite{goemans2002}. Related uses include approximation algorithms for weighted completion time on unrelated machines \cite{sitters2017} and completion-time scheduling with job assignments to machines \cite{sethuraman1999}. A job's completion time is the end of its processing. The objective here depends on all processing times, through an integral. Queyranne's fixed-machine result is attributed to his unpublished notes of 20 July 2010, identified in the proceedings note \cite[discussion~3 and ref.~4]{queyranne2010}. Those notes were unavailable for direct comparison.

J\"ager, Lindermayr and Megow \cite{jaeger2026} use a time-indexed weighted mean-busy-time relaxation to bound weighted completion time in polytope scheduling. Damerius, Kling and Schneider \cite{damerius2024} use fractional total completion time in a continuous LP to obtain a completion-time lower bound and an approximation guarantee for scheduling with a shared resource. Here the continuous first moment is itself the objective, and the algorithm recovers an exact optimum and an attaining schedule for identical machines.

\paragraph{Separated continuous linear programming.}
The fractional problem for a fixed horizon is an instance of the linear-data separated continuous linear program (SCLP) studied by Weiss \cite{weiss2008}. To see this, write $U_j(t)=\int_0^t u_j(s)\,\mathrm{d}s$ and impose the pointwise rate and total capacity constraints, introducing nonnegative slacks if needed. The cumulative rows can be written
\[
 U_j(t)\leq p_j,\qquad -U_j(t)\leq H-p_j-t .
\]
They follow from the terminal workloads and $0\leq u_j\leq1$, and they force $U_j(H)=p_j$. Maximizing $\int_0^H(H-t)\alpha\cdot u(t)\dt$ for these controls differs from minimizing the processing cost by the fixed term $H\alpha\cdot p$. Thus the continuous model already fits established machinery.

Weiss gives a simplex-based finite-step method. The author-posted revision bounds the number of pivots combinatorially under its feasibility and nondegeneracy hypotheses. A worst-case polynomial Turing bit bound would require polynomial control of both the number of steps and their bit costs in the variable-machine input size. Finite termination alone leaves both estimates open, and the literature examined did not supply that bound for this scheduling specialization. Shindin and Weiss extend simplex-type finite termination to a measure-valued continuous linear program with constant coefficients that includes SCLP \cite{shindinweiss2020}. Their general analysis allows exponential worst-case behavior. An implementation of a revised SCLP simplex is studied in \cite{shindin2021}. Its numerical results concern practical performance. Fleischer and Sethuraman give a polynomial approximation scheme with additive and multiplicative tolerances for a more general holding-cost formulation \cite{fleischer2005}. For any feasible half-line control, Tonelli's theorem gives $\int_0^\infty t u_j(t)\dt=\int_0^\infty(p_j-U_j(t))\dt$, which explains the holding-cost interpretation. The result here concerns exact rational recovery with polynomial worst-case bit time.

\paragraph{Purification and explicit realization.}
For a constant rate vector, decompose it into a distribution on the finite set of binary feasible job patterns. Classical finite-action atomless purification preserves the probability assigned to each pattern under both time measure and first-moment measure \cite[\S2]{dww1950}. On a signed interval use $(t-a)\dt$ for the second measure and recover the signed moment using the preserved work. Summing over the patterns containing a job preserves that job's work and first moment. This already gives measurable existence. The publisher's abstract of Anderson and Pullan \cite{andersonpullan1996} describes a constructive SCLP purification to an extreme point without decreasing the objective. The full text was unavailable for comparing its construction, precise hypotheses, and rational operation bounds. The realization proved here uses McNaughton's row wrap \cite[Theorem~3.1]{mcnaughton1959} on a half-interval and then reflects the operations. It gives explicit machine assignments and a finite rational output with a per-job record bound.

\paragraph{Optimization tools and scope of comparison.}
Parameter-dependent affine orderings and polyhedral descriptions are standard themes in parametric optimization \cite{nemesch2025}. Given the fixed-machine pattern formulation, a natural route is to dualize the $n$ workload constraints and seek exact recovery through the framework of Gr\"otschel, Lov\'asz and Schrijver \cite{gls1988}. With standard convex duality \cite{rockafellar1970}, the problem-specific steps are to describe the minimizers of an \emph{integrated} support function by polynomial-size rational rows, supply an exact workload-based separator, and turn the resulting rates into a moment-preserving schedule. The literature examined established prior continuous formulations, finite-step algorithms, and measurable purification. It did not identify the full variable-machine worst-case polynomial bit guarantee in \Cref{thm:q}. The unpublished notes, the uninspected portions of the purification literature, and the scope of the bibliographic coverage leave priority unresolved.
